\documentclass[10pt,a4paper]{article}

\usepackage[T1]{fontenc}
\usepackage[utf8]{inputenc}
\usepackage{lmodern}
\usepackage{amsmath,amssymb,amsthm}
\usepackage{booktabs,array,tabularx}
\usepackage{geometry}
\usepackage{xcolor}
\usepackage{enumitem}
\usepackage{microtype}
\usepackage{hyperref}
\usepackage{fancyhdr}

\geometry{margin=1in}
\hypersetup{
  colorlinks=true,
  linkcolor=blue!55!black,
  citecolor=blue!55!black,
  urlcolor=blue!55!black,
  pdftitle={Calving Iceberg: A Same-Message Context-Splitting Group-Key Recovery Attack},
  pdfauthor={Martin Feussner},
  pdfsubject={Cryptanalytic audit report for IACR ePrint 2026/1757},
  pdfkeywords={digital signatures, Iceberg, threshold signatures, MuSig2, nonce reuse, key recovery, BIP340}
}
\setlist{nosep,leftmargin=*}
\setlength{\parindent}{1em}
\setlength{\parskip}{0.25em}
\setlength{\emergencystretch}{2em}
\Urlmuskip=0mu plus 1mu
\pagestyle{fancy}
\fancyhf{}
\fancyhead[L]{Calving Iceberg}
\fancyhead[R]{18 October 2026}
\fancyfoot[C]{\thepage}
\renewcommand{\headrulewidth}{0.3pt}

\newtheorem{theorem}{Theorem}
\newtheorem{lemma}{Lemma}
\newcommand{\Zq}{\mathbb{Z}_q}
\newcommand{\KeyAgg}{\mathsf{KeyAgg}}
\newcommand{\PreRound}{\mathsf{PreRound}}
\newcommand{\SignRound}{\mathsf{SignRound}}
\newcommand{\Verify}{\mathsf{Verify}}

\title{\bfseries Calving Iceberg:\\
A Same-Message Context-Splitting Group-Key Recovery Attack}
\author{Martin Feussner\\
\small Selmer Center, University of Bergen\\
\small \texttt{martin.feussner@uib.no}}
\date{18 October 2026}

\begin{document}
\maketitle

\noindent\fcolorbox{black!55}{black!4}{\parbox{0.94\textwidth}{
This audit used the OpenAI model Daybreak Blue at the ultra reasoning level.
OpenAI Codex generated and selected this paper's title and deliberately chose
``Calving Iceberg'' as scheme-related wordplay: iceberg calving splits one ice
mass, while the attack splits one Iceberg session across several outer
contexts. OpenAI Codex performed AI-assisted target reconstruction,
attack-history searching, cryptanalysis, implementation, experimentation,
review coordination, and technical writing. An independently written
fresh-context AI checker reimplemented the secp256k1, BIP327, BIP340, sharing,
and extraction logic and returned PASS with negative controls. The exact
pinned-source C diagnostic was also rerun and returned PASS. The formal-game,
claim-boundary, and prior-art objections were reviewed separately before
publication. Independent human verification was not performed; it is optional
and is not a publication requirement.}}

\begin{abstract}
Iceberg, proposed in IACR ePrint 2026/1757, lets a threshold group occupy one
signer slot in an otherwise ordinary two-party MuSig2 execution. Its
deterministic group nonce pair is derived from a session identifier. The
paper's TS-MS-EUF-CMA game prevents two messages from being signed under one
identifier, but it consumes state separately for each honest member and does
not bind the adversary-chosen outer MuSig2 context across those members.

We give a practical full-parameter attack on the paper's concrete
$t=2,n=4$ secp256k1 profile. Corrupt one member. Query the other three members
once each under the same session identifier and the same message, while giving
each a different cosigner nonce. Combining each honest response with the
corrupt member's local response yields three complete group partial signatures
of the form
\[
 S_j=\varepsilon_j r_1+\varepsilon_j b_1b_{0,j}r_2
      +e_j a\gamma x.
\]
The coefficients are public and three ordinary contexts give a full-rank
$3\times3$ system. Solving it recovers the two fixed nonce scalars and the
group signing key $x$. The recovered key, together with the adversary's outer
cosigner key, produces a standard 64-byte BIP340 signature on a never-queried
message with no honest group participation.

An exact-source C diagnostic pinned to artifact commit
\texttt{b2f87cb047467dd78852bc32722a5228a7167f5e} recovered the key and verified
the fresh signature in 0.01 seconds using 2,560 KiB peak resident memory. An
independently written standard-library Python implementation reproduced the
result in 0.92 seconds using 15,360 KiB. Both use the ordinary BIP340 verifier,
and negative controls reject repeated contexts, two equations, omitted nonce
parity, and raw unaggregated responses. The three-response MuSig2 solve,
same-message context manipulation, and replicated-sharing split ingredients
are prior public knowledge; we claim only their concrete application to
Iceberg's formal game. No previous public report of this complete attack
against Iceberg was located in searches through 18 October 2026.
\end{abstract}

\section{Introduction}

Iceberg is a nested threshold multisignature intended to replace one side of a
Lightning channel's two-party MuSig2 signer with a threshold group while
leaving the counterparty and the on-chain BIP340 verifier unchanged
\cite{iceberg,bip327,bip340}. A key requirement is that the group commit to its
nonce before the transaction message and final signing quorum are known. To
make that possible, Iceberg derives two group nonce scalars deterministically
from replicated secret shares and a session identifier.

Deterministic MuSig nonces require more than message equality. A malicious
cosigner can keep the message fixed and change its nonce, thereby changing the
aggregate nonce coefficient and Schnorr challenge. This general phenomenon is
known from deterministic multisignature analyses, and the direct recovery of a
BIP327 key from three responses using one fixed two-scalar nonce pair is also
public \cite{musigdn,kohenreuse}. Iceberg attempts to prevent dangerous reuse
by requiring group agreement on one message per session identifier.

The formal game enforces exactly that message condition, but no corresponding
condition on the complete second-round input. Each honest member has its own
ready/done state. Consequently, three different honest members can each answer
once for one common message and identifier while seeing three different outer
cosigner nonces. Threshold interpolation turns these three distinct member
responses into three complete responses for the same group-level key and nonce
unknowns. The published three-response linear solve then recovers the group
key.

\paragraph{Contributions.}
We (i) give the exact game-admissible query schedule; (ii) derive the complete
BIP327 equation including both nonce and aggregate-key parity signs; (iii)
execute key recovery and a fresh BIP340 forgery against the pinned
full-parameter source; (iv) reproduce the result independently with a
separate secp256k1 implementation and negative controls; and (v) identify the
precise stronger assumption used by the proof but absent from Figure~3.

\paragraph{Originality boundary.}
We do not claim a new nonce-reuse technique. Kohen gives the three-response
BIP327 linear system and discusses replicated-sharing quorum splits
\cite{kohenreuse}; MuSig-DN discusses same-message context manipulation for a
deterministic one-nonce multisignature \cite{musigdn}; and the public Iceberg
source itself contains a commented observation that changing the
counterparty's nonce can change the signature hash without changing the
message. Our result is the concrete synthesis for Iceberg: its message-only
game binding and member-local state supply three legal complete group
responses in the claimed $2$-of-$4$ profile.

\section{Target and exact profile}

The target is the initial/current version of \emph{Enabling Threshold Custody
for the Lightning Network with Nested Threshold Multi-Signatures}, IACR ePrint
2026/1757, received 21 August and approved 22 August 2026
\cite{iceberg}. The retained 36-page PDF has SHA-256 digest
\begin{center}
\small\ttfamily
6a3bd9d2079551000ea4d69e18463f0975e85decbb5a37f0efff30e1c2508403.
\end{center}
The independently downloaded arXiv v1 source archive has SHA-256 digest
\begin{center}
\small\ttfamily
984b33ba56392a9e5aef5b64a0bd58c4946fe458dae3afd26452c8dfd1198ef5.
\end{center}

The retained profile is
\begin{center}
\small\ttfamily
ICEBERG-SECP256K1-T2-N4-MU3-NU2-DEPTH2-\\
BIP327-BIP340-SID32-ARTIFACT-B2F87CB.
\end{center}
It uses the prime-order secp256k1 group, threshold $t=2$, four members, one
static corruption, the paper's $\mu=2t-1=3$ live-member policy, two MuSig2
nonce components, 32-byte messages and identifiers, one fixed attacker-owned
outer cosigner, the group key first in the two-key list, and no BIP341 tweak.
This is the paper's smallest and benchmarked deployment profile, not a reduced
security parameter.

The author artifact is pinned at commit
\texttt{b2f87cb047467dd78852bc32722a5228a7167f5e} \cite{artifact}. The C
diagnostic includes its vendored secp256k1 implementation directly and calls
the native Iceberg share generation, public-share aggregation, nonce
generation, nonce aggregation, partial-signature generation and verification,
partial aggregation, BIP327 session, and BIP340 verification paths. No
implementation behavior is needed to create the weakness: the query schedule
is expressly permitted by the paper's game and the equations are those of the
construction.

\section{Iceberg signing equations}

Let $q$ be the secp256k1 group order and write scalar arithmetic in $\Zq$.
The threshold group holds signing scalar $x$ and public key $X=xG$. For a
session identifier $\mathit{sid}$, Iceberg deterministically derives two
shared nonce scalars $r_1,r_2$ and public points $R_1=r_1G$ and $R_2=r_2G$.
Its inner coefficient is
\[
 b_1=H_{\mathrm{Iceberg/noncecoef}}(R_1,R_2,X),
\]
using the source's extended point encodings. The group exports the ordinary
MuSig2 public nonce $(R_1,b_1R_2)$.

Fix an outer key list $[X,X_B]$, where the adversary owns $x_B$ and
$X_B=x_BG$. BIP327 key aggregation assigns public coefficient $a$ to $X$ and
forms the aggregate key $\widetilde X=aX+X_B$. Let $\gamma\in\{1,-1\}$ be
the source's aggregate-key parity correction. In outer context $j$, the
adversary supplies a fresh valid cosigner nonce pair. Aggregating it with the
fixed group nonce yields the public outer nonce coefficient $b_{0,j}$, final
nonce parity correction $\varepsilon_j\in\{1,-1\}$, and BIP340 challenge
$e_j$. Every one of these values is public.

If member $i$ holds Shamir-form shares $x_i,r_{1,i},r_{2,i}$, its second-round
response in context $j$ is
\begin{equation}\label{eq:member}
 s_{i,j}=\varepsilon_j r_{1,i}
 +\varepsilon_j b_1b_{0,j}r_{2,i}
 +e_j a\gamma x_i \pmod q.
\end{equation}
Lagrange interpolation of any $t$ valid member responses therefore produces
the group's ordinary MuSig2 partial scalar
\begin{equation}\label{eq:group}
 S_j=\varepsilon_j r_1
 +\varepsilon_j b_1b_{0,j}r_2
 +e_j a\gamma x \pmod q.
\end{equation}

For $t=2$, corrupt member 1 and honest member $i$ have interpolation weights
\[
 \lambda_1=\frac{i}{i-1},\qquad
 \lambda_i=-\frac{1}{i-1},
\]
so the adversary computes
$S_j=\lambda_1s_{1,j}+\lambda_is_{i,j}$. It knows $s_{1,j}$ from the corrupt
member's shares and obtains $s_{i,j}$ from one honest oracle call.

\section{Same-message context-splitting attack}

\subsection{Legal oracle schedule}

Figure~3 defines a global message map $M[\mathit{sid}]$ and a separate state
$\mathit{SID}[i][\mathit{sid}]$ for each member. The pre-round oracle changes
one member's state from unset to ready. The signing oracle checks only that
the selected member is ready and
\[
 M[\mathit{sid}]\in\{\bot,m\},
\]
then calls $\SignRound$ on the adversary-supplied aggregate pre-round message
and outer context, marks only that member done, and records the member in
$Q[m,\mathit{sid}]$. The paper explicitly describes the outer context as
including the group position, cosigner public keys, and outer aggregate
first-round message, and says that the adversary chooses it.

The following schedule is therefore admitted for one corrupt member:
\begin{enumerate}
  \item choose one session identifier $s$ and query $\PreRound(s,i)$ once for
        each honest member $i\in\{2,3,4\}$;
  \item aggregate those three valid pre-round contributions into the fixed
        group public nonce;
  \item fix one message $m$, one outer key list $[X,X_B]$, and one attacker
        cosigner key $x_B$;
  \item choose three distinct valid cosigner nonce pairs, producing contexts
        $\mathit{ctx}_2,\mathit{ctx}_3,\mathit{ctx}_4$;
  \item query honest member $i$ exactly once on
        $(s,i,m,\mathit{out},\mathit{ctx}_i)$.
\end{enumerate}
Every query uses the same $s$ and $m$, so the global message check succeeds.
Every member answers once, so its local state check succeeds. Only the
adversary-controlled cosigner nonce changes.

\subsection{Key extraction}

Combine each honest response with the locally generated corrupt response to
obtain $S_2,S_3,S_4$. Equation~\eqref{eq:group} gives
\[
\begin{pmatrix}S_2\\S_3\\S_4\end{pmatrix}
=
\begin{pmatrix}
\varepsilon_2 & \varepsilon_2b_1b_{0,2} & e_2a\gamma\\
\varepsilon_3 & \varepsilon_3b_1b_{0,3} & e_3a\gamma\\
\varepsilon_4 & \varepsilon_4b_1b_{0,4} & e_4a\gamma
\end{pmatrix}
\begin{pmatrix}r_1\\r_2\\x\end{pmatrix}
\pmod q.
\]
The adversary knows the matrix and right-hand side. Distinct outer nonces
change $b_{0,j}$ and $e_j$; ordinary choices give rank three except for a
negligible degeneracy that is detected before submitting the queries. The
deterministic reproduction selected candidates 0, 1, and 5. Their parity signs
were $(-1,-1,+1)$ and their matrix had rank three. Gaussian elimination
recovered $(r_1,r_2,x)$, and $xG$ exactly matched the public group key.

\subsection{Fresh-message forgery}

The adversary now knows both secret halves of the outer two-party signing
relation: recovered group scalar $x$ and its own cosigner scalar $x_B$. It
runs ordinary MuSig2 locally on a new message $m^*\ne m$ and outputs the
resulting 64-byte BIP340 signature under $\KeyAgg([X,X_B])$. The standard
BIP340 verifier accepts; changing one byte of $m^*$ makes it reject.

No honest member signs $m^*$. For every identifier,
$Q[m^*,\mathit{sid}]$ is empty. With $t=2$ and one corrupt member, Figure~3's
win condition is
\[
 |Q[m^*,\mathit{sid}]|=0<t-|C|=1.
\]
The result is therefore both secret-key recovery with fresh signing and an
accepted fresh-message forgery under the paper's claimed security game.

\begin{theorem}
In the Iceberg TS-MS-EUF-CMA game instantiated at $t=2,n=4$, a static
adversary corrupting one member recovers the group signing key after three
honest second-round queries on one message and one session identifier, and
then forges a fresh BIP340 signature with no honest participation, except with
the negligible probability that no tested context triple has full rank.
\end{theorem}

\begin{proof}
The three queries are legal by the member-local state and message-only global
check above. Each corrupt-plus-honest pair interpolates Equation~\eqref{eq:member}
to Equation~\eqref{eq:group}. A full-rank coefficient matrix uniquely reveals
$r_1,r_2,x$. Equality $xG=X$ publicly validates the recovered key. Knowledge
of $x$ and the controlled cosigner key permits ordinary BIP327 signing on a
fresh message, whose final signature is accepted by BIP340. The forged message
has no honest signing response, so it satisfies the game's freshness
inequality.
\end{proof}

\section{Proof gap and design status}

The first reduction says that the consensus assertion permits only one
message per identifier and concludes that one underlying signing-oracle call
per session suffices. The main reduction later assumes the stronger statement
that state agreement makes it impossible to query the same identifier on
``two different inputs to the second round.'' Figure~3 enforces equality only
of $m$; its other second-round inputs remain adversary-controlled. The three
queries above are therefore distinct second-round inputs but the same allowed
message.

This is a design-level attack on the formal construction, not an accidental
parser or implementation defect. It uses valid group elements, a claimed
parameter set, fewer than $t$ corruptions, authenticated member identities,
one response per honest member, one message per identifier, and the ordinary
outer verifier. It does not require rollback, repeated use by one member,
malformed encodings, weak keys, a hash collision, or an implementation race.
The pinned source execution is conformance evidence for the coefficients and
verifier endpoint, rather than the source of the weakness.

The attack assumes the adversarial outer signer and network powers already
granted by the TS-MS-EUF-CMA game. A deployment that authenticates and reaches
group-wide agreement on the entire second-round tuple before any member
answers would block this schedule, but that requirement is stronger than the
message agreement encoded in Figure~3 and used in the stated security claim.

\section{Experimental validation}

\subsection{Pinned-source diagnostic}

The exact-source C diagnostic used artifact commit
\texttt{b2f87cb047467dd78852bc32722a5228a7167f5e}. The retained attack source has
SHA-256
\begin{center}
\small\ttfamily
e9787f333a1cd09411629eae9d1751de5c21e29f5b61e8d13102c708ec624e01.
\end{center}
It was compiled with GCC 13.3.0 at C99 and optimization level \texttt{-O2}.
The program generated a trusted-dealer $t=2,n=4$ instance, used honest members
2, 3, and 4 for the three pre-round contributions, generated one response per
honest member in three contexts, verified every member response, completed and
verified all three queried-message outer signatures, solved the system,
regenerated the group public key, and verified a fresh-message signature.

The attack executable took 0.01 seconds, used 2,560 KiB maximum resident
memory, and used one worker. The local rerun reproduced every recorded scalar
and result flag. The only byte difference in captured standard output was an
extra trailing blank line in the sender's archived transcript.

\subsection{Independent reproduction}

A separately written 487-line Python checker uses only the standard library.
It implements secp256k1 point arithmetic, tagged hashes, Iceberg dealer and
VPSS derivation, Lagrange interpolation, BIP327 key and nonce coefficients,
BIP340 challenge, signing, and verification. Its adversary receives only the
public transcript, the corrupt member's shares, three honest responses, and
the controlled cosigner secret. Fixture-wide secrets are used only for
post-recovery assertions.

The checker returned PASS in 0.92 seconds with 15,360 KiB maximum resident
memory on one process and worker. Its JSON output matched the supplied
independent result byte for byte. It confirmed the following negative
controls:
\begin{itemize}
  \item two equations have rank two and do not determine three unknowns;
  \item repeating one outer context leaves the three-row matrix at rank two;
  \item omitting final-nonce parity makes the true solution violate the
        equations; and
  \item treating raw honest responses as complete group partials produces a
        scalar whose public point is not the group public key.
\end{itemize}

These are deterministic full-instance reproductions, not random-key
statistical trials. The algebra is generic, while the reported run establishes
one exact full-parameter key and transcript. The public package includes the
independent Python replay and the exact C diagnostic source.

\section{Prior work and attack-history search}

Kohen's March 2026 note gives the direct three-response MuSig2 solve for the
two fixed nonce scalars and signing key, and separately describes the danger of
splitting replicated nonce material among different quorums
\cite{kohenreuse}. Ruffing and Kohen also discuss the different limited-reuse
setting in which each BIP327 nonce pair is exposed only twice and many pairs
are combined with a generalized-birthday method \cite{ruffing}. The present
attack needs one nonce pair and three complete responses.

MuSig-DN describes same-message context manipulation in a deterministic
multisignature: a malicious cosigner changes its nonce and hence the challenge
while honest nonce material repeats \cite{musigdn}. Its displayed extraction
has one effective nonce scalar and does not use different threshold members.
The Iceberg authors' public explanation discusses a replicated-sharing split
across different messages and presents message agreement as the defense
\cite{icebergblog}. The attack here keeps the message fixed and varies only
the outer context that Figure~3 leaves unbound.

Searches through 18 October 2026 covered the exact paper title and identifier,
Iceberg and nested-MuSig aliases, authors, the artifact repository, issues,
pull requests, releases, scholarly and general search results, same-message
context attacks, MuSig2 nonce reuse, related Nested MuSig2 work, and the public
author material. No previous public report of the complete Iceberg-specific
attack was located. This is a dated, bounded statement; it does not exclude
private, deleted, unindexed, or later work, and it is not a claim that the
underlying techniques are new.

\section{Repair direction and limitations}

A repair must bind every honest member under one identifier to one complete
second-round computation, including the group aggregate pre-round message,
group position, outer key list, outer aggregate nonce, message, and any tweak
state. Checking the message alone is insufficient. Possible designs include a
group-agreed digest of the whole context or a revised state machine that stores
that digest before returning any share. Such a change must coexist with the
paper's requirement that the group nonce be fixed before the message exists,
and it needs a new proof; we do not claim a validated drop-in repair.

The result is limited as follows:
\begin{itemize}
  \item the executed source-backed endpoint is the claimed $t=2,n=4$,
        $\mu=3$, $\nu=2$ secp256k1 profile;
  \item the attack targets unforgeability in the paper's formal game and does
        not evaluate Lightning consensus, availability, privacy, or channel
        state safety beyond the stated interface;
  \item the source artifact is a benchmark prototype with a trusted dealer;
        the attack does not require or claim a DKG or production integration;
  \item the prior-art conclusion is bounded to public material located by the
        search date; and
  \item no novelty is claimed for the linear solve, deterministic-nonce
        context manipulation, or replicated-sharing quorum-split ideas.
\end{itemize}

\section*{Acknowledgements}

The computations were performed on the Norwegian Research and Education Cloud (NREC), using resources provided by the University of Bergen and the University of Oslo.

\begin{thebibliography}{9}

\bibitem{iceberg}
P. Gerhart, N. Kohen, J. Posner, and M. Furszyfer.
\newblock Enabling Threshold Custody for the Lightning Network with Nested
Threshold Multi-Signatures.
\newblock IACR Cryptology ePrint Archive, Report 2026/1757, 2026.
\newblock \url{https://eprint.iacr.org/2026/1757}.

\bibitem{artifact}
M. Furszyfer et al.
\newblock Iceberg benchmark artifact, commit
\texttt{b2f87cb047467dd78852bc32722a5228a7167f5e}.
\newblock \url{https://github.com/furszy/benchmark-iceberg/commit/b2f87cb047467dd78852bc32722a5228a7167f5e}.

\bibitem{bip327}
J. Nick, T. Ruffing, and E. J. Jin.
\newblock BIP 327: MuSig2 for BIP340-compatible Multi-Signatures, 2022.
\newblock \url{https://github.com/bitcoin/bips/blob/master/bip-0327.mediawiki}.

\bibitem{bip340}
J. Nick, T. Ruffing, and P. Wuille.
\newblock BIP 340: Schnorr Signatures for secp256k1, 2020.
\newblock \url{https://github.com/bitcoin/bips/blob/master/bip-0340.mediawiki}.

\bibitem{kohenreuse}
N. Kohen.
\newblock A Practical Attack Against Limited MuSig2 Nonce Reuse, 19 March
2026.
\newblock \url{https://nkohen.github.io/blog/musig2-nonce-reuse/}.

\bibitem{ruffing}
T. Ruffing.
\newblock How many nonce reuse before exposing your MuSig2 private key?,
Delving Bitcoin, 7 March 2024.
\newblock \url{https://delvingbitcoin.org/t/how-many-nonce-reuse-before-exposing-your-musig2-private-key/217/4}.

\bibitem{musigdn}
J. Nick, T. Ruffing, Y. Seurin, and P. Wuille.
\newblock MuSig-DN: Schnorr Multi-Signatures with Verifiably Deterministic
Nonces.
\newblock IACR Cryptology ePrint Archive, Report 2020/1057, 2020.
\newblock \url{https://eprint.iacr.org/2020/1057}.

\bibitem{icebergblog}
N. Kohen.
\newblock Iceberg: Nested Threshold Signatures for Lightning, 28 August 2026.
\newblock \url{https://nkohen.github.io/blog/iceberg/}.

\end{thebibliography}

\end{document}
